\chapter{Memory Management}
\label{chap:memmag}
\vspace{10.5cm}

\section{Memory Management}
Allocation means to assign, allot, distribute or "set apart for 
a particular purpose". Programs manage their memory paritioning or 
dividing it into different units performing specific tasks.
\newline
In C++, there are two ways to create objects and this aspect reflects on 
the existence of two different units used to allocate the memory,
the \textbf{stack} and \textbf{heap}, which manage the program's 
used/unused memory in two different ways. 

\begin{figure}[H]
    \begin{center}
        \includegraphics[scale=0.5]{OOPandUML/functional_memory_organization}
        \label{fig:functional_memory_organization}
        \caption{Functional memory organization of a running program}
    \end{center}
\end{figure}

\toclessnparagraph{Text}
This area contains machine instructions (i.e. the executable 
code). \textbf{Global} and \textbf{static} variables are typically 
allocated at program startup. \textbf{Stack} contains all the automatic 
variables. \textbf{Heap} is meant to allocate the memory for the dynamic 
variables through the \textbf{\texttt{new}} and \textbf{\texttt{delete}} 
operators. The space in between is allocated by the operating system but 
not used, and is meant to let both the heap and stack grow.

\toclessnparagraph{Stack}
This is a simple LIFO (last in, first out) data structure which must 
support at least the \textbf{push} and \textbf{pop} operations. Even 
though a stack only allows accesses at the top, it is easy to implement 
and fast to execute push and pop operations.

\toclessnparagraph{Heap}
This is a more flexible data structure, which allows programs to 
allocate or deallocate memory anywhere and in any convenient order. This 
might cause hole in the heap, i.e. unallocated memory surrounded by 
allocated memory.
\newline
The only restriction for the heap is that memory can be allocated in 
a contiguous block of memory large enough to satisfy the request with a 
single chunk of memory. This restriction has two effects: it requires 
the operating system to scan the heap until a contiguous block is found, 
and to coalesce adjacent blocks of memory returned by the program.
\newline
Side effect is that the operation on the heap are slower than the 
operations on the stack.

\subsection{Memory Allocation and Deallocation in C}
\subsubsection{\texttt{\href{https://en.cppreference.com/w/c/memory/malloc}{malloc}}}
\subsubsection{\texttt{\href{https://en.cppreference.com/w/c/memory/free}{free}}}

\subsection{Memory Allocation and Deallocation in C++}
\subsubsection{\texttt{\href{https://en.cppreference.com/w/cpp/language/new}{new expression}}}
\subsubsection{\texttt{\href{https://en.cppreference.com/w/cpp/language/delete}{delete expression}}}



